\subsection{Acquisition at more than one resolution}
The minimization over orders in the graph classification is performed at a specified predictive resolution. An acquisition rule intended to be reused at several resolutions has a different quantifier. Its observed history, rather than the requested accuracy, determines subsequent acquisitions. We prove that an entire curve produced by such a rule admits a single deterministic comparison order.

More precisely, call an acquisition rule resolution blind when its transition functions and independent seed law are fixed across predictive budgets, and it receives no feedback from a resolution-specific encoder. It may nevertheless read its whole revealed history. Section~\ref{sec:v27-policy} proves the simultaneous statement
\[
 \forall a\ \forall\sigma\ \exists\pi=\pi(a,\sigma)\ \forall M\ge1:
 \qquad R^{*,\mathrm{max}}_{\pi,M}(a)
                      \le C R^{\mathrm{o,av}}_{\sigma,M}(a).
\]
The constant is uniform in the calibration, the rule and the budget for the fixed experiment. The right-hand side even grants a full-history scheduler and a finite-prefix decoder, whereas the deterministic construction on the left uses the original charged memory. The order is selected before the predictive budget or codebook is specified. The proof selects a positive-mass trace under the first-block acquisition event and applies the acquired-coordinate small-ball bounds to that same trace at every budget. Conditional densities on selected histories are not required.

This statement leads to a characterization of finite reusable menus. A menu of at most $K$ randomized, history-dependent acquisition rules, with its index selected before data are observed, is matched to within bounded additive bits by at most $K$ deterministic orders. The optimum excess over a set of resolutions is consequently the finite order-set functional of \eqref{eq:v27-menu-functional}. Along analytic collision paths its leading coefficient is \eqref{eq:v27-leading-menu}. In the three-edge experiment, the interval penalty previously obtained for deterministic orders therefore holds for every single resolution-blind rule in the specified larger class. Two deterministic configurations remove that leading penalty.

The relevant information conventions are summarized below. The full-history scheduler is an explicitly stated converse relaxation; it is never used as an uncharged implementation of the upper bound.

\begin{center}
\small
\begin{tabular}{@{}p{.23\textwidth}p{.34\textwidth}p{.35\textwidth}@{}}
\toprule
Result & Rule across resolutions & Decoder information \\
\midrule
Pointwise graph law, Theorem~\ref{thm:v25-adaptive} & May be chosen anew for each budget & Label, query descriptor, and optionally the finite order prefix \\
Separator certificate, Theorem~\ref{thm:v26-separator} & No cross-resolution requirement & Label and visited-set query descriptor, not an extra order prefix \\
Whole-curve domination, Theorem~\ref{thm:v27-curve} & One fixed acquisition law; predictive encoders may change & The larger finite-prefix decoder is allowed in the converse \\
Menu characterization, Theorem~\ref{thm:v27-menu} & At most $K$ fixed conditional acquisition laws; selector precedes the tapes & A finite configuration index, plus the preceding oracle-decoder information \\
\bottomrule
\end{tabular}
\end{center}

The restrictions on reuse concern acquisition, not one common quantizer at every resolution. A rule that reads a budget-dependent compressed state can change its acquisition law when the budget changes and is not silently included in the whole-curve assertion.

Finally, Section~\ref{sec:v27-evaluated} supplies a fully evaluated statistical instance. For a positive two-trial detector, we derive the complete mixed law of the acquired posterior mean, reduce its finite-label average risk exactly to scalar quantization, and compute its sharp high-resolution coefficient. In the same experiment we evaluate the event mass, weighted coordinate density and query recovery norm in the separator bound. Removing the completion-word restriction doubles that complete certificate, but the certificate remains strictly below the sharp coefficient. This separates the quantitative value of the refinement from an assertion of sharp constants for the general graph theorem.
